% GENERATED FILE. Edit canonical lessons, then run npm run generate:books.
% canonical-source-hash: fnv1a64:0c1886b4285249b5
% canonical-lessons: FA-C11-cheshm, FA-C11-dast, FA-C11-pa, FA-C11-zaban

\chapter{Eye, Hand, Foot, Tongue}
\label{ch:fa-body}
\hlchaptermodality{11}

\begin{chapteropening}
\textbf{By the end of this chapter:} \emph{I can name my eye, hand, foot, and tongue in Persian, and I can say zabân for both tongue and language — the word for the thing this whole track teaches.}

The chapter and the tranche close on zabân, tongue and language together, then run a short scene built from every chapter so far: the name exchange, the wellbeing question, this chapter's own requests and family and body words, and khodâhâfez to close.
\end{chapteropening}

\section[cheshm]{\fa{چشم} — eye, and what \textquotedblleft{}how are you\textquotedblright{} was really pointing at}

\label{lesson:FA-C11-cheshm}



\begin{quote}
Say \textbf{\fa{دختر}}. Then read \textbf{\fa{حال} \fa{شما} \fa{چطور} \fa{است؟}} once more from the right edge — an earlier lesson asked about your \textbf{\fa{حال}}, your state, but never once named the parts of you that state belongs to. This chapter does.
\end{quote}

\subsection*{You'll want to know first — one word}
\begin{quote}
\textbf{\fa{چشم}} — \emph{cheshm} — \textbf{eye}
\end{quote}

Three letters, all already yours: \textbf{\fa{چ}} \emph{ch}, from \textbf{\fa{چیست}}; \textbf{\fa{ش}} \emph{sh}, from \textbf{\fa{شما}}; \textbf{\fa{م}} \emph{m}.

\begin{cousinweb}
\textbf{\fa{چشم}} continues Middle Persian \textbf{čašm}, tracing back through Old Iranian to the Indo-European root *\emph{spek'-}, “to look at, to observe.” English never inherited this root directly for its everyday word \textbf{eye} — that word comes from a different ancient root entirely — but it kept *\emph{spek'-} through Latin \emph{specere}, “to look at”: \textbf{spy}, \textbf{spectacle}, \textbf{species}, \textbf{inspect}, and \textbf{suspect} are all this same root, arriving sideways rather than head-on.
\end{cousinweb}

\subsection*{Your turn}
\begin{itemize}
\raggedright
  \item \emph{Say it:} \textbf{cheshm} — eye
  \item \emph{Connect:} \textbf{cheshm} $\leftarrow$ *\emph{spek'-} $\to$ Latin \textbf{specere} $\to$ English \textbf{spy}, \textbf{spectacle}, \textbf{species}
  \item \emph{Retrieve:} \textbf{\fa{حال} \fa{شما} \fa{چطور} \fa{است؟}} — the question this chapter's body words could actually answer
  \item \emph{Run through:} \textbf{mâdar, pedar, barâdar, dokhtar, cheshm} — five family and body words in one breath
\end{itemize}

\begin{tcolorbox}[breakable,colback=teal!4,colframe=teal!35!black,title={Before you move on}]
What does \textbf{\fa{چشم}} mean? (\textbf{Eye.}) Which three English words share its root, none of them meaning “eye”? (\textbf{Spy, spectacle, species.}) Which earlier question could a body-part chapter actually answer? (\textbf{\fa{حال} \fa{شما} \fa{چطور} \fa{است؟}})

Source: \href{https://en.wiktionary.org/wiki/\%DA\%86\%D8\%B4\%D9\%85}{Wiktionary: \fa{چشم}}.
\end{tcolorbox}
\section[dast]{\fa{دست} — hand, and a sound law that hides a cousin in plain sight}

\label{lesson:FA-C11-dast}



\begin{quote}
Say \textbf{\fa{چشم}} (\emph{cheshm}). This lesson's word looks nothing like its closest Indo-Iranian relative — and that mismatch is not an accident.
\end{quote}

\subsection*{You'll want to know first — one word}
\begin{quote}
\textbf{\fa{دست}} — \emph{dast} — \textbf{hand}
\end{quote}

Three letters, all already yours: \textbf{\fa{د} \fa{س} \fa{ت}} — \emph{d}, \emph{s}, \emph{t}.

\begin{cousinweb}
\textbf{\fa{دست}} continues Old Persian \textbf{dastā-}, from Indo-European *\emph{ǵ\textsuperscript{h}es-to-}, “hand.” Sanskrit's word for hand, \textbf{hasta}, is the very same root — and so, further out, is Greek \textbf{kheír}, which English kept indirectly in \textbf{surgery} (literally “hand-work”) and \textbf{chiropractic} (“done by hand”).

Old Persian's \textbf{d}, Avestan's \textbf{z}, and Sanskrit's \textbf{h} look like three unrelated sounds, but they are the completely regular Indo-Iranian outcome of the same ancestral sound in each of the three branches — the kind of correspondence a false cousin never has, and a real one always does.
\end{cousinweb}

\subsection*{Your turn}
\begin{itemize}
\raggedright
  \item \emph{Say it:} \textbf{dast} — hand
  \item \emph{Connect:} \textbf{dast} $\leftarrow$ *\emph{ǵ\textsuperscript{h}es-to-} $\to$ Sanskrit \textbf{hasta}, Greek \textbf{kheír} $\to$ English \textbf{surgery}, \textbf{chiropractic}
  \item \emph{Contrast:} Persian \textbf{d}, Avestan \textbf{z}, Sanskrit \textbf{h} — three sounds, one regular law
  \item \emph{Run through:} \textbf{cheshm, dast} — eye, hand
\end{itemize}

\begin{tcolorbox}[breakable,colback=teal!4,colframe=teal!35!black,title={Before you move on}]
What does \textbf{\fa{دست}} mean? (\textbf{Hand.}) Which Sanskrit word is its direct cousin? (\textbf{hasta}.) Which two English words trace to the same root through Greek \textbf{kheír}? (\textbf{Surgery}, \textbf{chiropractic}.)

Source: \href{https://en.wiktionary.org/wiki/\%D8\%AF\%D8\%B3\%D8\%AA}{Wiktionary: \fa{دست}}.
\end{tcolorbox}
\section[pâ]{\fa{پا} — foot, the plainest cousin in the whole track}

\label{lesson:FA-C11-pa}



\begin{quote}
Say \textbf{\fa{چشم}} and \textbf{\fa{دست}}. One more body word closes this chapter, and for once its history needs no caution at all.
\end{quote}

\subsection*{You'll want to know first — one word}
\begin{quote}
\textbf{\fa{پا}} — \emph{pâ} — \textbf{foot}
\end{quote}

Two letters, both already yours: \textbf{\fa{پ}} \emph{p}, from \textbf{\fa{پدر}}; long \textbf{\fa{ا}} \emph{â}.

\begin{cousinweb}
\textbf{\fa{پا}} continues Old Persian \textbf{pād-}, from Indo-European *\emph{ped-}/*\emph{pod-}, “foot.” Unlike \textbf{\fa{دست}}, this cousin needs no sound-law explanation to see: English \textbf{foot}, Latin \textbf{pes}/\textbf{pedis} (\textbf{pedal}, \textbf{pedestrian}), Greek \textbf{pous}/\textbf{podos} (\textbf{podiatry}, \textbf{tripod}), and Sanskrit \textbf{pāda} all sit close enough to hear the family resemblance directly. Some Indo-European cousins take a page of sound changes to prove; this one barely needs a sentence.
\end{cousinweb}

\subsection*{Your turn}
\begin{itemize}
\raggedright
  \item \emph{Say it:} \textbf{pâ} — foot
  \item \emph{Connect:} \textbf{pâ} $\leftarrow$ *\emph{ped-}/*\emph{pod-} $\to$ English \textbf{foot}, Latin \textbf{pes}, Greek \textbf{pous}, Sanskrit \textbf{pāda}
  \item \emph{Run through:} \textbf{cheshm, dast, pâ} — eye, hand, foot
  \item \emph{Name:} which of this chapter's three cousins needed a sound law, and which one did not
\end{itemize}

\begin{tcolorbox}[breakable,colback=teal!4,colframe=teal!35!black,title={Before you move on}]
What does \textbf{\fa{پا}} mean? (\textbf{Foot.}) Name four language families that keep a recognisable version of this word. (\textbf{English, Latin, Greek, Sanskrit.}) Which earlier body word needed a sound law to prove its cousin? (\textbf{\fa{دست}}, hand.)

Source: \href{https://en.wiktionary.org/wiki/\%D9\%BE\%D8\%A7}{Wiktionary: \fa{پا}}.
\end{tcolorbox}
\section[zabân]{\fa{زبان} — tongue, and language, and everything this track has built}

\label{lesson:FA-C11-zaban}



\begin{quote}
Say \textbf{\fa{چشم،} \fa{دست،} \fa{پا}} — eye, hand, foot. One organ this list has left out names something bigger than itself.
\end{quote}

\subsection*{You'll want to know first — one word, one new letter}
\begin{quote}
\textbf{\fa{زبان}} — \emph{zabân} — \textbf{tongue}, and also \textbf{language}
\end{quote}

Four letters, three of them already yours — \textbf{\fa{ب} \fa{ا} \fa{ن}} — plus one this track has never needed before: \textbf{\fa{ز}} \emph{ze}, an ordinary letter from the inherited Arabic set, read simply as English \emph{z}. \textbf{\fa{ز}-\fa{ب}-\fa{ا}-\fa{ن}}.

\begin{cousinweb}
\textbf{\fa{زبان}} continues Middle Persian \textbf{zabān}, tracing to the ancient Indo-European word for “tongue.” English \textbf{tongue} is the same root, worn down at home. Latin took the same root and, by a well-documented but still debated \emph{d}-to-\emph{l} shift, turned it into \emph{lingua} — the ancestor of English \textbf{language} and \textbf{linguistics}. Persian, like English, never split the organ from the ability: \textbf{\fa{زبانِ} \fa{فارسی}} \emph{zabân-e fârsi}, “the Persian tongue,” is simply the Persian \textbf{language} — the same word this whole track has been teaching you to use.
\end{cousinweb}

\subsection*{Your turn — everything this track has taught}
One short scene, built entirely from independently learned pieces, chapter by chapter:

\begin{itemize}
\raggedright
  \item \emph{Read it:} \textbf{\fa{اسمِ} \fa{من} ... \fa{است}} — the very first sentence this track built
  \item \emph{Read it:} \textbf{\fa{شما}} against \textbf{\fa{تو}} — the register this track built next
  \item \emph{Read it:} \textbf{\fa{حال} \fa{شما} \fa{چطور} \fa{است؟}} — the careful question that followed
  \item \emph{Say it:} \textbf{\fa{خوبم}} — the reply that answered it
  \item \emph{Say it:} \textbf{\fa{آب،} \fa{لطفاً}} / \textbf{\fa{کلیدِ} \fa{من،} \fa{لطفاً}} — this chapter's own requests
  \item \emph{Name:} \textbf{\fa{مادر،} \fa{پدر،} \fa{برادر،} \fa{دختر،} \fa{چشم،} \fa{دست،} \fa{پا،} \fa{زبان}} — eight new words, three chapters
  \item \emph{Read it:} \textbf{\fa{خداحافظ}}, and recall what it leaves unsaid — no verb, just the two halves an earlier lesson joined
  \item \emph{Connect:} \textbf{zabân} $\leftarrow$ old “tongue” word $\to$ English \textbf{tongue}, Latin \emph{lingua} $\to$ \textbf{language}, \textbf{linguistics}
\end{itemize}

\begin{tcolorbox}[breakable,colback=teal!4,colframe=teal!35!black,title={Before you move on}]
What does \textbf{\fa{زبان}} mean, in both of its senses? (\textbf{Tongue}, and \textbf{language}.) Which new letter did it need? (\textbf{\fa{ز}}, \emph{ze}.) Which two English words trace to the same root by a disputed sound shift? (\textbf{Tongue} directly; \textbf{language} and \textbf{linguistics} through Latin's \emph{d}-to-\emph{l} turn.) What is \textbf{\fa{زبانِ} \fa{فارسی}}? (The Persian language — the one this whole track teaches.)

Source: \href{https://en.wiktionary.org/wiki/\%D8\%B2\%D8\%A8\%D8\%A7\%D9\%86}{Wiktionary: \fa{زبان}}.
\end{tcolorbox}
